\section{Ensayos sobre planta simulada}
\label{sec:ensayos}

Todo lo anterior son demostraciones y comprobaciones numéricas de enunciados
aislados. Esta sección reporta las campañas de simulación —tres sobre una planta
con dinámica de ruedas, chasis y contactos, un banco reducido de trayectorias
continuas y una planta mixta con tres cargas— y declara con precisión qué
contrasta cada una y qué no.

\subsection{Misiones completas con pérdidas de red}

\begin{observacion}[Tres realizaciones de diez AMR y dos cargas]
\label{obs:ensayo-misiones}
Se ejecutan tres misiones con diez AMR —ocho asignados y dos en reserva— y dos
cargas de $1{,}84\times1{,}28$~m y $2{,}04\times1{,}38$~m, con catálogo de
cardinalidad cuatro por carga. Los parámetros de masa, inercia, semivía, radio de
rueda, par y batería son heterogéneos y sintéticos.

\begin{center}\small
\begin{tabular}{@{}llr@{}}
\toprule
Caso & Configuración & Tiempo físico \\
\midrule
S01 & \textsf{cargo} de $44$~kg y empuje de $62$~kg & $95{,}60$~s \\
S02 & dos cargas en \textsf{cargo}, $44$ y $62$~kg & $86{,}40$~s \\
S03 & mixto, con $5\%$ de pérdidas en la red de motor y retransmisión
  & $97{,}36$~s \\
\bottomrule
\end{tabular}
\end{center}

El caso S03 es el pertinente para \cref{obs:perdidas}: introducir un $5\%$ de
pérdidas con retransmisión cuesta $1{,}8\%$ de tiempo respecto de S01. Es una
medida, no una cota, y sobre una sola semilla por caso.

En S02 ambas cargas comparten el paso central en sentidos opuestos y una recibe
autorización antes que la otra, que es el comportamiento del testigo de
\cref{obs:posicionamiento-almacen}. Los vídeos proceden de estados y fuerzas
integrados, no de interpolación de caminos. La evaluación de cartografía se
ejecuta \emph{en sombra}: no modifica la planificación ni los motores, de modo
que estos ensayos no contrastan \cref{th:mapa} ni \cref{co:tensado}.
\end{observacion}

\subsection{Cruce negociado sin regla de prioridad}

El ensayo siguiente es el que más directamente interroga a este trabajo, porque
resuelve un cruce \emph{sin} el mecanismo que aquí se propone.

\begin{observacion}[Dos cargas cruzándose en un paso obligatorio]
\label{obs:ensayo-cruce}
Ocho AMR heterogéneos movilizan dos cargas de $44$ y $62$~kg que atraviesan en
sentidos opuestos el mismo cuello, de $5{,}4$~m de ancho, sin rodeo exterior
posible. No hay ruta previa, ni carril asignado, ni permiso por prioridad: cada
$0{,}4$~s cada coalición genera propuestas de movimiento y se selecciona la
pareja.

\begin{center}\small
\begin{tabular}{@{}lrr@{}}
\toprule
Métrica & Paso nominal & Obstáculo revelado \\
\midrule
Poses completas alcanzadas & $2/2$ & $2/2$ \\
Tiempo físico [s] & $48{,}40$ & $57{,}28$ \\
Energía total [Wh] & $1{,}195$ & $1{,}402$ \\
Centros simultáneos en el paso [s] & $20{,}16$ & $14{,}56$ \\
Menor separación monitorizada [m] & $0{,}096$ & $0{,}097$ \\
Pico de utilización del cono & $0{,}112$ & $0{,}149$ \\
\bottomrule
\end{tabular}
\end{center}

Dos lecturas. La primera, favorable a las hipótesis de este trabajo: el pico de
utilización del cono de fricción es de $0{,}112$ y $0{,}149$, es decir que la
adherencia opera al $11$–$15\%$ de su límite. El régimen ensayado está lejos de
la saturación de \cref{obs:multiplicador-friccion}, y por tanto \emph{no}
contrasta la ley de degradación en la zona donde predice parálisis.

La segunda es incómoda y conviene no esquivarla: las dos coaliciones ocupan el
paso simultáneamente durante $20{,}16$~s. El mecanismo de testigo de este trabajo
lo habría impedido, concediendo acceso exclusivo por turnos.
\end{observacion}

\begin{observacion}[Un ensayo que cuestiona la exclusividad]
\label{obs:ensayo-cuestiona}
El ensayo P01 lo lleva más lejos. Con un paso de $4{,}50$~m de ancho y
$5{,}20$~m de longitud —verificado como única conexión entre salas mediante un
test de conectividad— las dos cargas \emph{intercambian su orden longitudinal
dentro del pasillo}:

\begin{center}\small
\begin{tabular}{@{}lr@{}}
\toprule
Magnitud & Resultado \\
\midrule
Poses completas, con permanencia & $2$ de $2$ \\
Fin de la ejecución física & $52{,}16$~s \\
Intercambio de orden longitudinal & $25{,}92$~s \\
Ambas envolventes dentro del pasillo & $26{,}40$~s \\
Paradas dentro del pasillo & $0$~s en ambas \\
Mínima separación entre coaliciones & $0{,}3119$~m \\
Épocas de negociación & $131$ \\
\bottomrule
\end{tabular}
\end{center}

Cero paradas y ningún bloqueo, con las dos envolventes dentro del pasillo durante
$26{,}40$~s. Eso acota \cref{obs:ciclo-despeje}, donde se concluye que
el turno óptimo dura unas seis veces el tiempo de despeje, y de ahí se justifica
el testigo exclusivo. Ese análisis \emph{presupone} que el corredor se usa por
turnos. Si la renegociación continua permite uso simultáneo, el coste de despeje
desaparece porque no hay conmutación que pagar, y la comparación correcta no es
entre ciclos sino entre dos arquitecturas distintas.

Lo que el ensayo no dice es igual de importante. Consume $131$ épocas de
negociación —una cada $0{,}4$~s— y un predictor de pose por coalición, frente al
testigo, que exige un mensaje y un desempate. No hay garantía de espera acotada
como la del testigo, ni resultado de ausencia de bloqueo: hay un escenario en el
que no se produjo. La conclusión honesta es que la exclusividad compra garantía y
cuesta caudal, y que este ensayo mide lo segundo sin refutar lo primero.
\end{observacion}

\subsection{El historial de detenciones, leído con los teoremas}

Los ensayos anteriores son los que salieron bien. El registro de reparación
conserva también los que no, y su interés está en que cada detención lleva un
diagnóstico que este trabajo puede nombrar \cite{mayorga2026m1}.

\begin{observacion}[Cada modo de fallo tiene su enunciado]
\label{obs:historial-fallos}
\begin{center}\small
\begin{tabular}{@{}lccl@{}}
\toprule
Diagnóstico del banco & Entregas & Tiempo [s] & Enunciado que lo explica \\
\midrule
Guardia angular excesiva & $0/2$ & $180{,}00$ & guarda leída sobre el rango,
  no sobre la saturación (\cref{obs:rigidez-alcance}) \\
Contacto antes de \emph{stiction} & $2/2$ & $110{,}00$ & reacción lateral
  estática ausente en el modelo viscoso (\cref{obs:disco-compartido}) \\
Orientación \textsf{cargo} incompatible & $0/2$ & $13{,}76$ & autoridad
  degradada por geometría de formación (\cref{obs:incompatible}) \\
Autoridad antes de \emph{stiction} & $1/2$ & $160{,}00$ & parálisis por
  presupuesto: el óptimo existe y no llega a tiempo (\cref{obs:paralisis}) \\
Sin respaldo verificado & $1/2$ & $74{,}72$ & compromiso sin testigo completo
  (\cref{def:testigo}) \\
\bottomrule
\end{tabular}
\end{center}

El registro advierte que las ecuaciones y guardas cambiaron entre filas, de
modo que no es una comparación causal ni una campaña con modelo congelado. Su
valor aquí es distinto: muestra que las detenciones del banco no fueron
arbitrarias, sino instancias de fenómenos que este documento predice —y que en
dos casos, la reacción estática y la guarda sobre el rango, fueron precisamente
los que obligaron a corregir enunciados propios.
\end{observacion}

\begin{observacion}[Lo que verifica el código y lo que no]
\label{obs:pruebas-software}
La fuente de entrega supera $27/27$ pruebas nombradas, entre ellas: la
identidad Smith–par sobre diez mil casos con error $2{,}3\times10^{-13}$; la
equivalencia de presupuestos de barrera sobre diez mil casos con error
$1{,}6\times10^{-14}$; min-sum frente a enumeración sobre diez mil catálogos;
la proyección frente a un optimizador independiente con brecha
$-3{,}2\times10^{-14}$; los conos de impulso y la disipación lateral sobre diez
mil casos; y una reacción lateral estática no nula de $0{,}024$~N\,s, que es
justamente la que \cref{obs:disco-compartido} exige.

Con pérdidas del $1\%$ en la mensajería y sin retransmisión la ejecución se
detuvo a $23{,}92$~s por falta de instantánea confirmada; con hasta siete
retransmisiones por enlace se registraron $42\,007$ pérdidas y otras tantas
retransmisiones, y se completaron $2/2$ entregas. Conviene reproducir la
salvedad del registro tal cual: las trazas con y sin retransmisión coinciden
\emph{exactamente} en poses, pares, energía y tiempo para la misma semilla, y
eso es consecuencia del tiempo lógico síncrono, no una demostración de
compensación de retardos físicos. Las diez mil pruebas de cada familia tampoco
se suman como si fuesen misiones. Son verificaciones de identidades y de
software; la validación experimental de la construcción sigue siendo la que
\cref{obs:ensayos-alcance} acota.
\end{observacion}

\subsection{Banco reducido de trayectorias continuas}

Los ensayos siguientes proceden de una planta 2D reducida con dos agentes
compuestos —una coalición en \textsf{cargo} de cuatro AMR bajo soportes
solidarios, con huella rígida y radio de giro derivado de sus ruedas, y un
empujador— y contrastan la construcción de \cref{sec:continuacion}
\cite{mayorga2026compendio}.

\begin{observacion}[Trayectoria emergente sin \emph{waypoint}]
\label{obs:ensayo-r6}
Los únicos datos de navegación son estado actual, objetivo, mapa de
obstáculos, huella y la última propuesta del otro agente; la recta
origen--destino es solo inicialización y las continuaciones se renegocian cada
$0{,}35$--$0{,}40$~s. En la corrida principal: tiempo de misión $27{,}30$~s,
longitudes $16{,}44$~m y $19{,}73$~m, ocupación máxima $3{,}752$ y separación
mínima $0{,}0309$~m, sin colisión. Un optimizador central con toda la
información que planifica antes de mover completó en los mismos $27{,}3$~s con
pequeñas diferencias de distancia y energía, lo que muestra que ese central
no es un oráculo certificado; el reactivo local cayó en bloqueo. La separación
de $3{,}1$~cm es la que motiva \cref{obs:reserva-robusta}.
\end{observacion}

\begin{observacion}[Treinta mundos, siete métodos]
\label{obs:ensayo-r10}
Treinta mundos pareados —diez cruces, diez pasillos anchos, diez radiales— con
la misma realización de origen, destino, obstáculo y parámetros para cada
método:

\begin{center}\small
\begin{tabular}{@{}lccccc@{}}
\toprule
Método & Éxito seguro & Cruce & Pasillo ancho & Radial & CPU mediana [s] \\
\midrule
Reactivo local & $0/30$ & $0/10$ & $0/10$ & $0/10$ & --- \\
ORCA reducido (uniciclo) & $10/30$ & $0/10$ & $10/10$ & $0/10$ & $0{,}132$ \\
Prioridad con primitivas (PIBT) & $10/30$ & $0/10$ & $10/10$ & $0/10$ & $0{,}351$ \\
MPC distribuido por intercambio & $28/30$ & $9/10$ & $10/10$ & $9/10$ & $11{,}56$ \\
Megajuego (horizonte deslizante y precio) & $27/30$ & $9/10$ & $10/10$ & $8/10$ & $11{,}09$ \\
Central NMPC con toda la información & $28/30$ & $10/10$ & $10/10$ & $8/10$ & --- \\
Mejor central (cota derivada) & $29/30$ & & & & \\
\bottomrule
\end{tabular}
\end{center}

En comparación pareada con corrección de Holm, el megajuego supera
significativamente al reactivo, a ORCA y a la prioridad en tasa de éxito; no
se detecta diferencia frente al MPC distribuido ni al central con $N=30$. En
los veintisiete mundos seguros comunes con el MPC distribuido, las diferencias
medianas son $0{,}000\%$ en tiempo, $+0{,}028\%$ en distancia y $+0{,}032\%$
en energía, ninguna significativa; frente al central, $+1{,}613\%$,
$-0{,}038\%$ y $+0{,}094\%$, con tiempo y energía detectables y distancia no.
Por escenario: en el cruce las medianas son $31{,}6$~s (MPC distribuido),
$30{,}8$~s (megajuego) y $29{,}8$~s (central), y es el régimen donde la
predicción es necesaria; en el pasillo ancho todos los no reactivos completan
y la prioridad con primitivas es la más rápida ($24{,}325$~s frente a
$25{,}0$~s); en el radial el MPC distribuido gana un caso al megajuego, que es
la evidencia de \cref{obs:homotopia}. Los dos métodos predictivos usaron unos
$524$ intercambios lógicos medianos: el banco no respalda una afirmación de
menor cómputo ni menor comunicación del megajuego.
\end{observacion}

\begin{observacion}[Alcance del banco reducido]
\label{obs:ensayo-r10-alcance}
La planta es 2D reducida, la energía es un proxy
$E=\int\bigl(n_{\mathrm{AMR}}P_0+c_vv^2+c_\omega\omega^2\bigr)\mathrm dt$, las
implementaciones de los métodos comparados son reducidas y no las de
referencia, y $N=30$ no resuelve diferencias pequeñas. El propio informe
propone la campaña confirmatoria: cincuenta a cien mundos por escenario con
semilla congelada, implementaciones de referencia, márgenes robustos de
$0{,}20$, $0{,}30$ y $0{,}50$~m, barrido de $\varepsilon=T_{\mathrm{game}}/T_{\mathrm{phys}}$,
revisión homotópica, y repetición sobre la planta con ruedas y contacto.
Hasta entonces, lo que el banco sostiene es la necesidad de predicción en los
escenarios difíciles y la paridad con el MPC distribuido en la capa continua.
\end{observacion}

\subsection{Tres cargas, dos modos físicos y un solo paso}

\begin{observacion}[Planta mixta con ruedas dinámicas]
\label{obs:ensayo-mixed2d}
Diez AMR heterogéneos —masas de $19$ a $36$~kg, parachoques de $0{,}47$ a
$0{,}59$~m, radios de rueda de $8{,}5$ a $11$~cm, pares de $4{,}2$ a $7$~N\,m
por motor— transportan tres cargas por una planta de $22\times14$~m con un
único paso de $3{,}30$~m de ancho: C1 en \textsf{cargo} ($60$~kg,
$1{,}60\times1{,}10$~m, tres AMR), C2 en \textsf{caging} por empuje perimetral
($45$~kg, octógono de $1{,}25\times0{,}95$~m, cuatro AMR) y C3 en
\textsf{cargo} ($95$~kg, $1{,}90\times1{,}30$~m, tres AMR)
\cite{mayorga2026compendio}. Cada rueda tiene velocidad angular, inercia, par
y reacción de tracción propios, con el disco de fricción compartido de
\cref{pr:disco-friccion}; los apoyos de \textsf{cargo} son pads con giro
libre y fuerza tangencial por deformación y velocidad relativa objetiva
(\cref{pr:pad-energia}); en \textsf{caging} la fuerza aparece solo al contacto
del parachoques, $f_n=[k(\delta)]_+$, y no hay tracción a distancia.

\begin{center}\small
\begin{tabular}{@{}lrrr@{}}
\toprule
Carga & Entra al paso [s] & Sale [s] & Entrega aceptada [s] \\
\midrule
C1 (\textsf{cargo}) & $19{,}64$ & $34{,}72$ & $63{,}04$ \\
C3 (\textsf{cargo}) & $51{,}16$ & $72{,}76$ & $101{,}40$ \\
C2 (\textsf{caging}) & $85{,}68$ & $104{,}88$ & $131{,}64$ \\
\bottomrule
\end{tabular}
\end{center}

El orden $1\to3\to2$ no está escrito en ningún controlador. La corrida termina
a $134{,}68$~s con las tres cargas dentro de tolerancias de $0{,}26$~m,
$0{,}065$~m/s y $0{,}12$~rad, sin una sola penetración no deseada; los mínimos
del integrador fueron $0{,}117$~m a paredes, $0{,}235$~m entre parachoques y
$0{,}830$~m entre cargas, con a lo sumo una envolvente en el paso. Las
deformaciones medianas de los pads fueron $0{,}79$ y $1{,}14$~mm con máximo
$9{,}66$~mm —elásticas, no deslizamiento—, y en \textsf{caging} hubo $3{,}80$
contactos activos de cuatro en promedio, con compresión de parachoques de
hasta $5{,}48$~mm. Cinco semillas nominales completaron $3/3$; el
refinamiento a $1$~ms desplazó la última entrega de $131{,}64$ a
$132{,}04$~s; y un caso de fricción superior muy baja salió del dominio de
adherencia y se detuvo a $0{,}16$~s, que es el límite declarado de
\cref{obs:disco-compartido}.
\end{observacion}

\begin{observacion}[Lo que cuesta organizar turnos, medido]
\label{obs:mixed2d-exclusividad}
La ablación sin exclusividad del paso, con la misma protección anticolisión
de referencias, entrega también $3/3$ y termina en $94{,}04$~s frente a
$134{,}68$~s, con $17{,}92$~s de dos envolventes dentro del recurso y cero
penetraciones. Es la tercera medida independiente del mismo hecho que
\cref{obs:ensayo-cruce} y \cref{obs:ensayo-cuestiona}: la capacidad unitaria
del corredor compra la garantía de \cref{th:ciclo-espera} y paga caudal, y
en este banco el precio es un $30\%$ del tiempo. No demuestra que los turnos
sean superiores para cualquier almacén ni que la ocupación compartida sea
segura en general; dice cuánto vale, en esta planta, la ausencia de
interbloqueo.
\end{observacion}

\begin{observacion}[Alcance de la planta mixta]
\label{obs:mixed2d-alcance}
Los equipos están fijados —\textsf{cargo} empieza montado y \textsf{caging}
próximo al objeto—, la vertical es cuasiestática, y la coordinación es una
referencia centralizada: mapa completo y un programa cuadrático cooperativo que
modifica referencias de velocidad. No es el cierre distribuido a par del
megajuego. El confinamiento en \textsf{caging} es observado, no certificado:
el certificado topológico de \cref{th:cerco} no se evaluó en línea. El
paquete supera ocho pruebas —conservación de momento del contacto, identidad
disipativa del pad, equilibrio vertical, límites de tracción y rango del mapa
de empujes—, y una prueba de rango detectó una geometría inicial de
\textsf{caging} defectuosa, que es exactamente el modo de fallo que
\cref{obs:incompatible} predice.
\end{observacion}

\begin{observacion}[Alcance de las campañas sobre la planta con ruedas]
\label{obs:ensayos-alcance}
Las salvedades que los propios informes de las tres campañas sobre la planta con ruedas declaran son estas. Los equipos
parten con contacto ya adquirido: no se ensaya reclutamiento, acoplamiento ni
descarga, de modo que \cref{th:potencial}, \cref{co:terminacion} y todo el
\cref{sec:standby} quedan sin contrastar. Se resuelve exactamente el juego finito
de propuestas de cada época con un predictor reducido, y no se afirma optimalidad
global de la misión no lineal. Los parámetros son sintéticos y no corresponden a
hardware identificado. En el ensayo con obstáculo, las propuestas de la carga en
\textsf{cargo} se restringen a traslación con orientación objetivo nula, lo que el
informe declara explícitamente. Y cada caso se ejecuta con una semilla.

Con ese alcance, lo que las campañas aportan es evidencia de que la planta cierra
—diez AMR, dos cargas, misiones completas, con pérdidas de red— y una medida del
régimen de adherencia en que opera. No son la validación de la construcción
teórica de este documento, y presentarlas como tal sería excederse.
\end{observacion}
